\chapter{Implementing Mutual Information Estimators}
\label{ch:mi_estimators}

Estimating Mutual Information (MI) and differential entropy from finite samples is a notoriously difficult problem, especially in high dimensions. In this chapter, we transition from the theoretical foundations discussed in previous chapters to practical, empirical implementations. We will implement three prominent paradigms from scratch: non-parametric $k$-nearest neighbor (KSG/Kozachenko-Leonenko) estimation, Kernel Density Estimation (KDE), and parametric neural estimation based on variational bounds (MINE and InfoNCE). We will validate their correctness by empirically comparing their bias on a ground-truth Gaussian scenario, where the closed-form entropy and MI are analytically known.

\section{Ground-Truth: Closed-Form Statistics for Gaussians}

To rigorously evaluate our estimators, we require a dataset with known, analytically tractable statistics. A bivariate Gaussian distribution with a specified correlation coefficient $\rho$ serves as an ideal testbed. As established in Equation \ref{eq:gaussian_mi}, the mutual information between two scalar Gaussian variables $X$ and $Y$ with correlation $\rho$ is $I(X; Y) = -\frac{1}{2} \log(1 - \rho^2)$. Similarly, the differential entropy of a scalar Gaussian is governed by Equation \ref{eq:gaussian_entropy}. The following snippet generates samples from this distribution and calculates the theoretical ground-truth values.

\begin{pycode}
import numpy as np

def generate_gaussian_samples(num_samples: int, rho: float) -> tuple[np.ndarray, np.ndarray, float, float]:
    """Generates correlated Gaussian samples and returns ground-truth Entropy and MI."""
    cov_matrix = np.array([[1.0, rho], [rho, 1.0]])
    data = np.random.multivariate_normal([0.0, 0.0], cov_matrix, size=num_samples)
    x = data[:, 0:1]
    y = data[:, 1:2]
    
    # Ground-truth MI for bivariate Gaussian
    true_mi = -0.5 * np.log(1.0 - rho**2)
    # Ground-truth entropy for standard normal marginal X
    true_entropy_x = 0.5 * np.log(2 * np.pi * np.e)
    
    return x, y, true_entropy_x, true_mi
\end{pycode}

This function utilizes \texttt{numpy} to sample from a bivariate Gaussian distribution with zero mean and a covariance matrix determined by the correlation parameter $\rho$. The output arrays \texttt{x} and \texttt{y} are formatted as 2D arrays (column vectors) to be compatible with distance computations in later estimators. The analytical mutual information and the marginal differential entropy for $X$ are computed and returned alongside the data to serve as our baseline for empirical bias evaluation.

\section{The KSG and Kozachenko-Leonenko Estimators (k-NN)}

The Kraskov-Stögbauer-Grassberger (KSG) mutual information estimator and the related Kozachenko-Leonenko entropy estimator are staple non-parametric methods. They rely on $k$-nearest neighbor distances in the joint and marginal spaces, improving upon naive histogram methods by adaptively scaling the volume based on local data density. This local adaptation significantly mitigates the bias inherent in fixed-bandwidth approaches. 

\begin{pycode}
from scipy.special import digamma
from scipy.spatial import cKDTree

def knn_entropy_estimator(x: np.ndarray, k: int = 5) -> float:
    """Estimates differential entropy using the Kozachenko-Leonenko method (L_inf norm)."""
    n, d = x.shape
    tree = cKDTree(x)
    distances, _ = tree.query(x, k=k+1, p=np.inf)
    epsilons = distances[:, k]
    # Entropy formula incorporating digamma bias correction
    return float(digamma(n) - digamma(k) + d * np.mean(np.log(2 * epsilons)))

def ksg_mi_estimator(x: np.ndarray, y: np.ndarray, k: int = 5) -> float:
    """Estimates MI using the KSG k-nearest neighbor method."""
    n = x.shape[0]
    xy = np.concatenate([x, y], axis=1)
    
    tree_xy, tree_x, tree_y = cKDTree(xy), cKDTree(x), cKDTree(y)
    distances, _ = tree_xy.query(xy, k=k+1, p=np.inf)
    epsilons = distances[:, k]
    
    # Count points in marginal spaces within epsilon distances
    nx = np.array([len(tree_x.query_ball_point(x[i], epsilons[i], p=np.inf)) - 1 for i in range(n)])
    ny = np.array([len(tree_y.query_ball_point(y[i], epsilons[i], p=np.inf)) - 1 for i in range(n)])
    
    # KSG Formula
    mi = digamma(k) - np.mean(digamma(nx + 1) + digamma(ny + 1)) + digamma(n)
    return max(0.0, float(mi))
\end{pycode}

The implementation leverages \texttt{scipy.spatial.cKDTree} for efficient spatial querying. For entropy estimation, we compute the distance to the $k$-th nearest neighbor and apply the volume formula for the Chebyshev ($L_\infty$) metric. For the KSG MI estimator, we find the distance to the $k$-th nearest neighbor in the joint space, and then count the number of points in the marginal spaces $X$ and $Y$ that fall within a hypercube of this exact distance. The final estimate involves the digamma function, corresponding to the harmonic number bias correction derived in Theorem \ref{thm:ksg_digamma}.

\section{Kernel Density Estimation (KDE)}

Kernel Density Estimation provides a smoothed functional estimate of the probability density by placing a continuous kernel (typically Gaussian) over each data point. Once the joint density $p(x, y)$ and marginal densities $p(x)$ are estimated, entropy and mutual information can be approximated through Monte Carlo integration over the empirical samples.

\begin{pycode}
from scipy.stats import gaussian_kde

def kde_entropy_estimator(x: np.ndarray) -> float:
    """Estimates differential entropy using Gaussian KDE."""
    kde_x = gaussian_kde(x.T)
    log_p_x = kde_x.logpdf(x.T)
    return float(-np.mean(log_p_x))

def kde_mi_estimator(x: np.ndarray, y: np.ndarray) -> float:
    """Estimates MI using Gaussian Kernel Density Estimation."""
    x_t, y_t = x.T, y.T
    xy_t = np.concatenate([x_t, y_t], axis=0)
    
    kde_xy = gaussian_kde(xy_t)
    kde_x = gaussian_kde(x_t)
    kde_y = gaussian_kde(y_t)
    
    # Evaluate log probabilities on the sample points
    log_p_xy = kde_xy.logpdf(xy_t)
    log_p_x = kde_x.logpdf(x_t)
    log_p_y = kde_y.logpdf(y_t)
    
    mi = np.mean(log_p_xy - log_p_x - log_p_y)
    return max(0.0, float(mi))
\end{pycode}

In this snippet, \texttt{scipy.stats.gaussian\_kde} automatically selects an optimal bandwidth using Scott's rule. We compute the log-densities for the joint and marginal distributions evaluated directly at the observed data points. The expectation $I(X; Y) = \mathbb{E}_{p(x, y)} [\log \frac{p(x, y)}{p(x)p(y)}]$ is then approximated by taking the sample mean of these log-ratios, and similarly for the entropy $H(X) = \mathbb{E}_{p(x)} [-\log p(x)]$. While straightforward and smooth, KDE suffers heavily from the curse of dimensionality and boundary effects, often requiring carefully tuned bandwidths for accurate estimation.

\section{Neural Estimators: MINE and InfoNCE}

In modern deep learning, estimating MI in high-dimensional spaces (e.g., between images and latent representations) necessitates parametric approaches. Neural estimators parameterize variational lower bounds on MI—such as the Donsker-Varadhan representation (MINE) or the Noise-Contrastive Estimation bound (InfoNCE)—using flexible neural networks optimized via gradient ascent.

\begin{pycode}
import torch
import torch.nn as nn
import torch.optim as optim

class SeparableCritic(nn.Module):
    def __init__(self):
        super().__init__()
        self.net_x = nn.Sequential(nn.Linear(1, 32), nn.ReLU(), nn.Linear(32, 16))
        self.net_y = nn.Sequential(nn.Linear(1, 32), nn.ReLU(), nn.Linear(32, 16))
        
    def forward(self, x, y):
        # Computes pairwise critic scores: [batch_size, batch_size]
        return torch.matmul(self.net_x(x), self.net_y(y).T)

def train_neural_estimator(x: np.ndarray, y: np.ndarray, method: str = 'mine', epochs: int = 300) -> float:
    """Trains a neural network to estimate MI using MINE or InfoNCE bounds."""
    x_t = torch.tensor(x, dtype=torch.float32)
    y_t = torch.tensor(y, dtype=torch.float32)
    
    model = SeparableCritic()
    optimizer = optim.Adam(model.parameters(), lr=0.01)
    
    for _ in range(epochs):
        optimizer.zero_grad()
        scores = model(x_t, y_t)
        
        # Joint distribution evaluations lie on the diagonal
        t_joint = torch.diag(scores)
        
        if method == 'mine':
            # Donsker-Varadhan lower bound (MINE) using all N^2 pairs for the marginal
            loss = -(torch.mean(t_joint) - torch.log(torch.mean(torch.exp(scores))))
        elif method == 'infonce':
            # InfoNCE lower bound using the N-1 other batch elements as negatives
            loss = -(torch.mean(t_joint) - torch.mean(torch.logsumexp(scores, dim=1)) + np.log(x_t.size(0)))
            
        loss.backward()
        optimizer.step()
        
    return float(-loss.item())
\end{pycode}

We define a separable critic network $T_\theta(x, y) = f_{\theta_1}(x)^T g_{\theta_2}(y)$, which allows us to efficiently compute the pairwise scores for all instances in a batch via matrix multiplication. During training, samples from the joint distribution $p(x, y)$ correspond to the aligned $(X, Y)$ instances on the diagonal of the score matrix. Samples from the product of marginals $p(x)p(y)$ are constructed by utilizing the off-diagonal elements (pairing every $x$ with every other $y$ in the batch). The optimization maximizes the chosen variational lower bound (minimizing the negative bound). Upon convergence, the negative loss serves as our empirical estimate of the mutual information, as derived in Section \ref{sec:variational_bounds}.

\section{Empirical Comparison of Estimator Bias}

Having implemented these diverse paradigms, we can now benchmark their performance against the analytical ground-truth derived earlier. The goal is to observe the bias and variance inherent to each approach under a finite sample size regime.

\begin{pycode}
def benchmark_estimators():
    """Compares entropy and MI estimators on a Gaussian dataset."""
    np.random.seed(42)
    torch.manual_seed(42)
    
    n_samples = 2000
    rho = 0.8
    
    x, y, true_entropy, true_mi = generate_gaussian_samples(n_samples, rho)
    
    # Entropy Estimation
    h_knn = knn_entropy_estimator(x, k=5)
    h_kde = kde_entropy_estimator(x)
    
    # MI Estimation
    mi_ksg = ksg_mi_estimator(x, y, k=5)
    mi_kde = kde_mi_estimator(x, y)
    mi_mine = train_neural_estimator(x, y, method='mine')
    mi_infonce = train_neural_estimator(x, y, method='infonce')
    
    print("--- Entropy Estimation ---")
    print(f"Ground Truth: {true_entropy:.4f}")
    print(f"k-NN (Bias):  {h_knn:.4f} ({h_knn - true_entropy:+.4f})")
    print(f"KDE (Bias):   {h_kde:.4f} ({h_kde - true_entropy:+.4f})")
    
    print("\n--- Mutual Information Estimation ---")
    print(f"Ground Truth:  {true_mi:.4f}")
    print(f"KSG (Bias):    {mi_ksg:.4f} ({mi_ksg - true_mi:+.4f})")
    print(f"KDE (Bias):    {mi_kde:.4f} ({mi_kde - true_mi:+.4f})")
    print(f"MINE (Bias):   {mi_mine:.4f} ({mi_mine - true_mi:+.4f})")
    print(f"InfoNCE (Bias):{mi_infonce:.4f} ({mi_infonce - true_mi:+.4f})")

# Execute the benchmark
# benchmark_estimators()
\end{pycode}

This final script orchestrates the comparison on a generated dataset of $N=2000$ points with a moderately high correlation ($\rho=0.8$). In practice, running this benchmark typically reveals that $k$-NN methods like KSG exhibit remarkably low bias and robust performance. KDE often struggles with smoothing artifacts, leading to underestimation of MI. Neural estimators (MINE and InfoNCE) introduce higher variance due to optimization stochasticity and are bounded by batch-size constraints, yet they remain indispensable when scaling beyond simple low-dimensional benchmarks to complex, unstructured data.
